\documentclass[11pt]{article}
\usepackage{fontspec}
\usepackage[french]{babel}
\usepackage{amsmath}
\usepackage{amssymb}
\usepackage{amsthm}
\usepackage{mathrsfs}
\usepackage{tikz-cd}
\usepackage{enumitem}

\newcommand{\printedpage}[1]{\marginpar{\footnotesize p.~#1}}
\newcommand{\msp}[1]{\marginpar{\footnotesize [#1]}}
% chapter line of the table: title (bold as printed) and page range
\newcommand{\tochead}[2]{\par\medskip\noindent #1\hfill #2\par\smallskip}
% entry of the table: indentation, label, text, page
\newcommand{\tocl}[4]{\par\noindent\hangindent=\dimexpr #1+3em\relax\hspace*{#1}\makebox[2.5em][r]{#2}\hspace{0.5em}#3\dotfill #4}
\newcommand{\Ch}{\underline{\mathrm{Ch}}}

\begin{document}

\printedpage{1}
% head: (none)

\begin{center}
{\LARGE\textbf{Les Dérivateurs}}

\bigskip

{\large Alexandre Grothendieck}
\end{center}

\bigskip

\tochead{\textbf{I. Généralités sur les (pré)dérivateurs.}}{[1]-[126]}
\tocl{0em}{1.}{Préliminaires sur les foncteurs images inverses des préfaisceaux sur des catégories d'indices, et leurs adjoints multiples}{[1]}
\tocl{0em}{2.}{Préliminaires heuristiques sur les dérivateurs}{[4]}
\tocl{0em}{}{Morphismes homotopiquement propres, homotopiquement lisses}{[53]}
\tocl{0em}{}{Notions de $\mathbb{D}$-équivalence et de $\mathbb{D}$-contractibilité}{[61]}
\tocl{0em}{}{$\mathbb{D}$-fibrations}{[76]}

\tochead{\textbf{II. Cofinalité (à droite et à gauche) (préliminaire à la cofinalité cohomologique).}}{[1]-[99]}
\tocl{0em}{1.}{Cofinalité cohomologique}{[1]}
\tocl{0em}{2.}{Les catégories localement filtrantes sont les catégories localement totalement 0-connexes}{[11]}
\tocl{0em}{3.}{Retour sur la cofinalité}{[37]}
\tocl{0em}{4.}{Foncteurs biasphériques}{[49]}
\tocl{0em}{5.}{Dualité formelle}{[69]}
\tocl{0em}{6.}{Transporteurs, noyaux, stabiliseurs relatifs à $W$}{[83]}
\tocl{0em}{7.}{Morphismes cofiltrants et localement cofiltrants}{[90]}

\tochead{\textbf{III. Hom externes dans les dérivateurs.}}{[1]-[12]}
\tocl{0em}{1.}{Relation des dérivateurs avec la catégorie Hot}{[1]}

\tochead{\textbf{IV. Diagrammes substantiels.}}{[73]-[83]}
\tocl{0em}{}{Construction de diagrammes substantiels pour un dérivateur}{[73]}

\tochead{\textbf{V. Catégories de chemins et localisation.}}{[1]-[52]}
\tocl{0em}{1.}{La catégorie des chemins}{[1]}
\tocl{0em}{2.}{Calcul d'une catégorie de fractions $X\Sigma^{-1}$}{[5]}
\tocl{0em}{3.}{Passage aux catégories de diagrammes}{[28]}

\tochead{\textbf{VI. HOT.}}{[1]-[204]}
\tocl{0em}{}{Le dérivateur HOT}{[1]}
\tocl{0em}{}{Nouvelle version de la théorie de $\mathrm{HOT}_W$}{[171]}

\clearpage
\printedpage{2}
% head: 2

\tochead{\textbf{VII. Catégories de chemins (2).}}{[1]-[171]}
\tocl{0em}{1.}{Intervalles et leur variantes}{[1]} % [sic] leur variantes -> leurs variantes, accord du pluriel
\tocl{0em}{2.}{Sous-catégories pleines d'un intervalle}{[10]}
\tocl{0em}{3.}{Chemins. Catégorie de chemins d'un type donné $\tau$}{[13]}
\tocl{0em}{4.}{Catégorie de chemins $\Ch(X)$}{[24]}
\tocl{0em}{5.}{Propriétés des morphismes d'intervalles}{[31]}
\tocl{0em}{6.}{Composition dans \underline{Ty}}{[48]}
\tocl{0em}{7.}{Explicitation des catégories de fractions $X\Sigma^{-1}$ en termes des $\Ch(X;a,b)$}{[59]}
\tocl{0em}{8.}{Le lemme d'homotopie pour les morphismes d'intervalles}{[83]}
\tocl{0em}{9.}{Théorème d'homotopie pour les catégories de chemins}{[94]}
\tocl{0em}{10.}{Question de comparaison pour les catégories Cat-fibrées satisfaisant au “théorème d'homotopie” (base filtrante)}{[99]}
\tocl{0em}{11.}{Rapide retour sur les $(\mathcal{M},\Sigma)$, dans l'optique des dérivateurs}{[129]}
\tocl{0em}{12.}{Retour sur les ‘catégories de chemins’ dans Cat, comme catégorie de modèles pour Hot}{[137]}
\tocl{0em}{13.}{Où on prouve enfin que $\Ch(X;-,-)$ est aussi bon (voire meilleur) que $\Ch_\infty(X;-,-)$}{[157]}

\tochead{\textbf{VIII. 1-types d'homotopie relatifs : leur intégration\,\dots}}{[1]-[9]}
\tocl{0em}{1.}{Intégration des 1-types d'homotopie relatifs sur une catégorie}{[1]}
\tocl{0em}{2.}{Questions de changement de base}{[6]}

\tochead{\textbf{IX. Retour sur les catégories $\mathcal{M}W^{-1}$, comparaison avec Quillen.}}{[1]-[25]}
\tocl{0em}{1.}{Fonctorialité de $\mathrm{Ch}^{\Sigma}_{\infty}(a,b)$ en $a$, $b$ dans $\mathcal{M}\Sigma^{-1}$}{[1]}
\tocl{0em}{2.}{Comparaison avec Quillen}{[14]}
\tocl{0em}{3.}{Essai de construction du foncteur $\Omega$}{[24]}

\tochead{\textbf{X. Comparaison de Cat avec $\Delta^\wedge$.}}{[1]-[5]}
\tocl{0em}{1.}{Rappel des notions de $\mathbb{D}$-asphéricité, et lien avec le ‘localisateur fondamental $W_{\mathbb{D}}$’}{[1]}

\tochead{\textbf{XI. Hot-fibrations, foncteurs propres et foncteurs lisses etc.}}{[127]-[136]}
\noindent\hspace*{1.5em}(dans Cat).\par\smallskip
\tocl{0em}{1.}{Hot-fibrations dans Cat}{[127]}

\tochead{\textbf{XII. Caractérisation de $W_\infty$. Foncteurs $W$-propres, $W$-lisses etc. Sommes amalgamées et carrés $W$-cocartésiens dans Cat.}}{[1]-[275]}
\tocl{0em}{I}{Caractérisation de $W_\infty$}{[1]}
\tocl{0em}{II}{Perplexités sur les axiomes et la terminologie}{[18]}
\tocl{0em}{III}{Classes remarquables de foncteurs liés à $W$}{[22]}
\tocl{0em}{IV}{Sur l'axiome W(4) et ses variantes}{[53]}
\tocl{0em}{V}{Sommes amalgamées et carrés $W$-cocartésiens dans Cat}{[59]}
\tocl{2em}{1.}{Généralités sur $U\overset{i}{\hookrightarrow}X\overset{j}{\hookleftarrow}T$ et les sommes amalgamées}{[59]}
\tocl{4.5em}{1.1.}{Données de recollement entre un ouvert et un fermé}{[59]}
\tocl{4.5em}{1.2.}{Relation avec le recollement des topos}{[61]}

\clearpage
\printedpage{3}
% head: 3

\tocl{4.5em}{1.3.}{Flèches dans Cat, de source ou but un objet recollé $(X,U,T)$}{[65]}
\tocl{4.5em}{1.4.}{Construction de sommes amalgamées à partir d'une immersion ouverte où fermée}{[66]} % [sic] où fermée -> ou fermée, conjonction (cf. 2.2)
\tocl{2em}{2.}{Sommes amalgamées et $W$-équivalences}{[75]}
\tocl{4.5em}{2.1.}{Axiome du carré cocartésien et axiome du cube (ou ‘Lemme des cinq’)}{[75]}
\tocl{4.5em}{2.2.}{Application aux carrés cocartésiens déduits d'une immersion ouverte ou fermée}{[80]}
\tocl{4.5em}{2.3.}{Objets totalement $W$-asphériques de Cat, et morphismes totalement $W$-asphériques}{[87]}
\tocl{4.5em}{2.4.}{Application aux immersions ouvertes dwyériennes}{[107]}
\tocl{4.5em}{2.5.}{Complément : critère d'asphéricité totale d'une catégorie $W$-propre sur une autre}{[115]}
\tocl{2em}{3.}{Sommes amalgamées et intégration}{[121]}
\tocl{4.5em}{3.1.}{Réduction canonique au cas d'un diagramme d'immersions fermées}{[121]}
\tocl{4.5em}{3.2.}{Comparaison de l'intégrale avec une somme amalgamée ordinaire. Carrés $W$-homotopiquement cocartésiens}{[128]}
\tocl{4.5em}{3.3.}{Intégration propre versus intégration lisse. Énoncé d'équivalence entre les deux points de vue. Commentaires sur le cas des limites $W$-homotopiques quelconques}{[133]}
\tocl{4.5em}{3.4.}{Critères pour les carrés homotopiquement cocartésiens}{[146]}
\tocl{4.5em}{3.5.}{Applications aux sommes amalgamées dans une catégorie $X^\wedge$ (avec $X\in\mathrm{Ob}\,\mathrm{Cat}$)}{[158]}
\tocl{0em}{VI}{Critère de $W_\infty$-asphéricité pour les $\Omega^n(X,x)$}{[171]}
\tocl{2em}{1.}{Cogitations sur $W_n$}{[171]}
\tocl{2em}{2.}{Critère de $W_n$-asphéricité}{[187]}
\tocl{2em}{3.}{Sur la philosophie des fibres homotopiques}{[191]}
\tocl{2em}{4.}{Axiomatisation des ‘fibres $W$-homotopiques’}{[207]}
\tocl{2em}{5.}{Axiomatique des $W$-types d'homotopie relatifs non localement constants $\mathrm{Ho}^{\mathcal{L}}_{W}(S)$}{[218]}
\tocl{2em}{6.}{Cofibre $W$-homotopique dans Cat}{[227]}
\tocl{2em}{7.}{Retour sur l'axiomatique de la fibre $W$-homotopique}{[240]}
\tocl{2em}{8.}{Appendice : Retour sur les axiomes W(7), W(7 bis) (pages 75, 76)}{[265]}

\tochead{\textbf{XIII. Catégories de modèles (1).}}{[1]-[98]}
\tocl{0em}{1.}{Couples de Quillen $(\Phi,\Psi)$}{[1]}
\tocl{0em}{2.}{Catégories de Quillen faibles, et construction d'icelles (Théorème-scholie conjectural)}{[7]}
\tocl{0em}{3.}{Catégories de Quillen (deuxième mouture)}{[27]}
\tocl{0em}{4.}{Les variances $u^*$, $u_*$, $u_!$ pour $(\mathcal{M}(I),W(I))$ (I variable)}{[44]}
\tocl{0em}{5.}{Catégories de modèles quillénisables}{[50]}
\tocl{0em}{6.}{Catégories de modèles quillénisables, et foncteurs $u_!$, $u_*$}{[57]}
\tocl{0em}{7.}{Filtrations cardinales et parties accessibles d'une grosse catégorie}{[77]}

\clearpage
\printedpage{4}
% head: 4

\tochead{\textbf{XIV. Carrés h-cartésiens et h-cocartésiens.}}{[1]-[8]}
\tocl{0em}{1.}{Axiomes de structures quadrangulées (une tentative avortée)}{[1]}
\tocl{0em}{2.}{Exemples}{[7]}

\tochead{\textbf{XV. Théorèmes de factorisation. (Modèles (2)).}}{[1]-[137]}
\tocl{0em}{1.}{Construction fondamentale : la suite transfinie de factorisation d'une flèche $f$}{[1]}
\tocl{0em}{2.}{Variation des factorisations $f=p_\alpha i_\alpha$ avec la partie $C_0$ génératrice}{[12]}
\tocl{0em}{3.}{Majoration des cardinaux}{[16]}
\tocl{0em}{4.}{Compléments sur les structures de Quillen}{[28]}
\tocl{0em}{5.}{Résultats téchniques préliminaires sur les parties accessibles de $\mathrm{Ob}\,\mathcal{M}$, stables par limites inductives filtrantes}{[43]} % [sic] téchniques -> techniques, accent
\tocl{0em}{6.}{Essai avorté pour un critère d'engendrement des parties L-accessibles de $\mathrm{Fl}(\mathcal{M})$}{[73]}
\tocl{0em}{7.}{Construction de triples de Quillen clos avec $C=\mathrm{Mon}\subset\mathrm{Cof}_W$}{[84]}
\tocl{0em}{8.}{Application aux catégories de modèles $(\widehat{A},W_A)$}{[100]}
\tocl{0em}{9.}{Révision des notations, pour les axiomes d'un localiseur fondamental}{[116]}
\tocl{0em}{10.}{Application au dérivateur $\mathrm{HOT}_W$}{[134]}

\tochead{\textbf{XVI. Localiseurs fondamentaux dans Cat.}}{[1]-[124]}
\tocl{0em}{1.}{Axiomes Loc\,(1) à Loc\,(3) (et variantes): localiseurs fondamentaux}{[1]}
\tocl{0em}{2.}{Axiomes Loc\,(4) et Loc\,(5) (des sommes directes, de connexité)}{[8]}
\tocl{0em}{3.}{Axiome Loc\,(6) du carré cocartésien}{[16$'$]}
\tocl{0em}{4.}{Axiome Loc\,(7) des $W$-fibrations}{[27]}
\tocl{0em}{5.}{Axiome des limites Loc\,(8), et l'axiome Loc\,(9) d'accessibilité}{[71]}
\tocl{0em}{6.}{La condition Loc\,$(5,\mathbb{D}_{W_\omega})$ implique $W\subset\mathrm{W}_\infty$}{[83]}
\tocl{0em}{7.}{Commentaires sur les axiomes}{[97]}

\tochead{\textbf{XVII. Catégories à fibrations et à cofibrations. (Modèles (3)).}}{[1]-[118]}
\tocl{0em}{1.}{Préliminaires}{[1]}
\tocl{0em}{2.}{Description de $C_I\subset\mathrm{Fl}(\mathcal{M}(I^\circ))$, pour $I$ ensemble ordonné}{[25]}
\tocl{0em}{3.}{Catégories de modèles à cofibrations}{[39]}
\tocl{0em}{4.}{Catégories à cofibrations spéciales}{[71]}
\tocl{0em}{5.}{Relation des catégories à cofibrations spéciales avec les catégories de Quillen spéciales}{[84]}
\tocl{0em}{6.}{Application à la construction des dérivateurs (associés aux triples de Quillen spéciaux)}{[98]}

\tochead{\textbf{XVIII. Catégories et ensembles accessibles.}}{[50]-[492]}
\tocl{0em}{4.}{Paratopos et l'opération $\mathcal{M}\boxtimes\mathcal{N}$}{[50]}
\tocl{2em}{4.1.}{Pseudo-topos. Problème du produit}{[50]}
\tocl{2em}{4.2.}{Cafouillages sur $\mathcal{M}\boxtimes\mathcal{N}$}{[57]}
\tocl{2em}{4.3.}{L'opération $\mathcal{M}\boxtimes\mathcal{N}$ et l'axiome d'accessibilité}{[64]}
\tocl{2em}{4.4.}{Étude de $\mathrm{Ind}_\pi(C)$ (pour $C$ petite catégorie stable par $\varinjlim$ de cardinal $\leq\pi$)}{[79]}

\clearpage
\printedpage{5}
% head: 5

\tocl{2em}{4.5.}{Filtration cardinale canonique d'une catégorie $\mathcal{U}$-paratopos}{[106]}
\tocl{2em}{4.6.}{Caractérisation des catégories paratopos (démonstration de 4.3.6)}{[139]}
\tocl{2em}{4.7.}{$\mathrm{Ind}_\pi C$ ($C$ quelconque) et filtrations cardinales généralisées}{[161]}
\tocl{2em}{4.8.}{Catégories $\mathcal{U}$-accessibles}{[187]}
\tocl{2em}{4.9.}{Le théorème de représentation indicielle pour $\mathrm{Ind}_\pi(C)$}{[217]}
\tocl{2em}{4.10.}{Le théorème de représentation indicielle : démonstration}{[243]}
\tocl{2em}{4.11.}{Application à l'existence des limites dans une catégorie de la forme $\mathrm{Ind}_{\pi_0}(C)$}{[267]}
\tocl{2em}{4.12.}{Foncteurs accessibles entre catégories accessibles}{[273]}
\tocl{2em}{4.13.}{Ensembles accessibles d'objets d'une catégorie accessible}{[281]}
\tocl{2em}{4.14.}{Critères d'accessibilité pour grosses catégories}{[370]}
\tocl{2em}{4.15.}{Catégorie accessible sur une autre}{[425]}
\tocl{2em}{4.16.}{Exemples et questions}{[464]}
\tocl{2em}{4.17.}{Retour sur les paratopos, et sur les opérations $\mathcal{M}\boxtimes\mathcal{N}$, $\underline{\mathrm{Hom}}_!(\mathcal{M},\mathcal{N})$}{[492]}
\par\noindent\hspace*{1.5em}\textbf{Appendice}.\hfill[1]-[19]
\tocl{2em}{}{Ultraproduits de corps, et catégories accessibles}{[1]}
\tocl{2em}{}{Application à un exemple d'une sous-catégorie non accessible}{[11]}

\tochead{\textbf{XIX. Modèles (4).}}{[1]-[50]}
\tocl{0em}{1.}{Rappels sur les couples de Quillen (couples de Quillen spéciaux)}{[1]}
\tocl{0em}{2.}{Rappels sur les triples de Quillen clos et les triples de Quillen spéciaux}{[13]}
\tocl{0em}{3.}{Cas des catégories à cofibrations, des catégories à fibrations}{[33]}

\bigskip
\noindent\rule{4cm}{0.4pt}

{\footnotesize\noindent\hangindent=2em \textsc{Note}. Dans cette table de matières, ainsi que dans l'édition électronique provisoire des “Dérivateurs”, on suit le découpage de l'ouvrage effectué a posteriori par Grothendieck. Les numéros des pages sont ceux de l'original, indiquées entre crochets dans le texte numérisé. Le chapitre XI, qui commence à la page 127, était à l'origine la dernière section du chapitre I, qui se termine à la page 126. De même, le chapitre XVIII était la section 4 du chapitre XIX.\par}

\end{document}
